\subsection{Measured savings, compatibility overhead and whole recognition}

The maintained suite passes all 1,090 tests with Regina in 218.217 seconds.
The 12 existing batch tests and six new integration tests take 0.556 and
0.045 seconds respectively. The final source-pinned audit takes 8.101 seconds.
Three separate benchmark executions contain 180 kernel calls, 120 group-stage
calls and 380 whole-recognition calls, plus 36, 24 and 76 retained warm-ups.
All 680 measured calls complete. Each fixture has five rounds with shuffled
old/new and old/new A/A arms. All four runs preserve the same 424 native Python
and corpus hashes; the pinned old package is loaded separately. Kernel,
stage and pipeline runs take 39.082, 2.843 and 34.649 seconds respectively.

Kernel timing includes fresh arena construction and complete checked raw
substitution. The shared family has $k$ independent definitions
$g_i=(ab)^{2^{4096}}$ sharing one source power grammar, with both signs used
in retained roots. The tower family has $k$ dependent definitions with
noncommuting contexts, each raised to $2^{128}$. The legacy family uses the
same towers but reverses their witness order, deliberately exercising the
compatibility fallback. No represented long word is expanded. Outputs have
matching pinned-old and current structural digests; literal correctness is
separately checked by the small-system audit.

\begin{center}\small
\begin{tabular}{@{}lrrrrrr@{}}
\toprule
Family & pivots & old ms & new ms & old/new & old A/A & new A/A \\
\midrule
shared & 8 & 100.833 & 44.068 & 2.266 & 0.988 & 0.996 \\
shared & 32 & 267.890 & 44.679 & 6.007 & 1.000 & 0.994 \\
shared & 128 & 948.063 & 45.743 & 20.552 & 0.993 & 0.984 \\
tower & 8 & 14.265 & 11.313 & 1.258 & 0.977 & 1.003 \\
tower & 32 & 60.742 & 47.640 & 1.284 & 1.008 & 0.988 \\
tower & 128 & 272.603 & 215.662 & 1.243 & 0.978 & 1.046 \\
legacy & 8 & 14.085 & 15.152 & 0.924 & 0.992 & 0.999 \\
legacy & 32 & 60.553 & 67.022 & 0.903 & 0.991 & 0.958 \\
legacy & 128 & 276.179 & 308.654 & 0.895 & 0.953 & 0.990 \\
\bottomrule
\end{tabular}
\end{center}


At 128 shared pivots, old/new total medians are 948.063/45.743 ms, with
paired ratio 20.552. The replay work counter falls from 2,699,059 to 55,213.
Exactly 4,357 source nodes are summarized once; the represented retained
word length has 4,099 bits. Both implementations add 4,100 nodes in this
case, so the saving is repeated inspection of shared grammars, not output
size. The 128-level tower instead demonstrates virtual inverse savings:
new allocations fall from 50,044 to 33,407, and replay work from 391,056 to
221,082. The paired total-time ratio is 1.243, with 16,514-bit output lengths.

The legacy compatibility cost is visible. On reversed 128-level towers,
old/new medians are 276.179/308.654 ms and the paired ratio is 0.895. New
allocations remain 50,044; replay work rises from 391,056 to 442,138 because
the charged rank scan precedes the unchanged fallback. The smaller reversed
cases also regress, with ratios 0.924 and 0.903. Kernel old A/A medians range
from 0.953 to 1.008 and new A/A from 0.958 to 1.046. The large shared-case gain
and the legacy overhead both remain recorded. Newly emitted ordered witnesses
avoid this fallback; compatibility does not imply identical finite-budget cost.

Group-stage timing uses validated stabilized-circle diagrams. It includes
fresh diagram construction, recovery, the same optional batch portfolio and
mandatory independent source replay, while bypassing earlier simplification.

\begin{center}\small
\begin{tabular}{@{}lrrrrr@{}}
\toprule
Crossings & old ms & new ms & old/new & old A/A & new A/A \\
\midrule
8 & 1.430 & 1.274 & 1.154 & 1.032 & 1.000 \\
16 & 2.800 & 2.479 & 1.139 & 0.996 & 1.011 \\
32 & 5.695 & 4.816 & 1.180 & 0.993 & 0.997 \\
64 & 11.500 & 9.748 & 1.188 & 0.995 & 1.000 \\
128 & 24.006 & 20.978 & 1.213 & 1.012 & 1.032 \\
256 & 47.905 & 51.442 & 1.128 & 0.977 & 0.997 \\
\bottomrule
\end{tabular}
\end{center}


Paired old/new medians range from 1.128 to 1.213. The largest case shows why
paired ratios and separate arm medians must be distinguished: at 256 crossings,
old/new medians are 47.905/51.442 ms, but the five paired ratios are
0.931, 1.128, 1.213, 1.084 and 1.208. The last two rounds also have substantially
larger absolute times. Median old/new A/A controls across this family range
from 0.977 to 1.032 and 0.997 to 1.032, but these do not capture all individual
variation. These data suggest modest stage savings, not a uniform factor.

A separate deterministic replay audit checks the saved stage proofs under
both verifiers, 24 times in total. At 256 crossings, compressed arena work
falls from 31,010 to 18,097 and final arena nodes from 3,049 to 2,038. That
counter includes SLP construction and replay after independent diagram
reconstruction, but excludes the preceding source-recovery budget. Producer
work rises from 94,320 to 95,842 to construct the ordering; the producer still
has 2,544 nodes and the certificate remains 12,775 bytes. These counters
explain why local verification savings give a smaller whole-stage effect.
The additional driver and records are \path{data/ordered_batch_replay_work.py}
and \path{data/ordered-batch-replay-work.json}.

Whole-recognition timing uses the original 19-diagram corpus and the same
optional batch portfolio on both versions. It includes fresh validated inputs,
all earlier stages, bounded search and mandatory certificate replay. No
incomplete result is omitted; all 380 measured calls complete.

\begin{center}\small
\begin{tabular}{@{}lrrrrr@{}}
\toprule
Input & old ms & new ms & old/new & old A/A & new A/A \\
\midrule
survivor-00 & 14.267 & 14.144 & 1.003 & 1.004 & 0.998 \\
survivor-01 & 27.925 & 28.734 & 0.968 & 1.003 & 0.998 \\
survivor-02 & 9.355 & 9.249 & 1.010 & 0.983 & 0.990 \\
survivor-03 & 38.030 & 38.350 & 0.990 & 1.002 & 1.010 \\
mirror-03 & 5.924 & 5.816 & 1.033 & 1.017 & 1.008 \\
survivor-04 & 28.923 & 28.983 & 1.000 & 1.000 & 1.001 \\
survivor-05 & 15.295 & 15.443 & 0.983 & 1.004 & 1.007 \\
survivor-06 & 5.278 & 5.151 & 1.028 & 1.003 & 0.990 \\
survivor-07 & 19.346 & 19.521 & 0.993 & 1.003 & 1.004 \\
survivor-08 & 7.067 & 6.978 & 0.998 & 0.993 & 0.979 \\
mirror-08 & 15.225 & 15.164 & 0.999 & 0.990 & 0.987 \\
survivor-09 & 5.586 & 5.464 & 1.016 & 1.006 & 1.007 \\
survivor-10 & 5.179 & 5.158 & 1.004 & 1.004 & 1.000 \\
survivor-11 & 4.765 & 4.701 & 1.013 & 1.001 & 0.999 \\
gordian & 1207.402 & 1210.837 & 0.998 & 1.001 & 0.992 \\
trefoil & 0.058 & 0.056 & 1.038 & 1.010 & 0.980 \\
figure\_eight & 0.062 & 0.062 & 1.000 & 0.992 & 0.981 \\
conway & 0.974 & 0.975 & 1.002 & 1.001 & 1.003 \\
kinoshita\_terasaka & 1.002 & 1.010 & 0.998 & 1.000 & 0.999 \\
\bottomrule
\end{tabular}
\end{center}


The ratios range only from 0.968 to 1.038; old A/A medians range from 0.983 to
1.017 and new A/A from 0.979 to 1.010. Gordian medians are 1,207.402 and
1,210.837 ms, with paired ratio 0.998. This workload demonstrates no useful
whole-recognizer acceleration from the new checker. Earlier exits and the
existing discovery fallback limit its contribution. The strong kernel result
therefore concerns compressed ordered proof replay; it is not promoted to a
recognition speedup or a general complexity result.
